\begin{center}
    ВЫВОД
\end{center}

В ходе выполнения лабораторной работы было продолжено изучение процессора с регистровой архитектурой, 
освоены форматы команд и данных. 

В соответствии с индивидуальным заданием было выполнено формирование адресов числовых значений, 
их перевод в шестнадцатеричную систему счисления и последующая запись в память. 

На основании заданного алгоритма была составлена и отлажена программа арифметико-логической 
обработки целочисленных данных с использованием различных способов косвенной адресации. 
Проведенный контроль результатов путем сравнения данных ручной трассировки с результатами выполнения программы 
в симуляторе подтвердил правильность работы разработанного программного кода и корректность использования механизмов адресации операндов.